% Generated draft: requires mathematical and source review.
\documentclass[11pt]{article}
\usepackage[margin=25mm]{geometry}
\usepackage{fontspec}
\setmainfont{Latin Modern Roman}
\usepackage{amsmath,amssymb,mathtools,unicode-math}
\setmathfont{Latin Modern Math}
\usepackage{enumitem,needspace,xurl,hyperref}
\setlength{\parindent}{0pt}
\setlength{\parskip}{5pt}
\setlength{\emergencystretch}{3em}
\newcommand{\N}{\mathbb{N}}
\newcommand{\Z}{\mathbb{Z}}
\newcommand{\Q}{\mathbb{Q}}
\newcommand{\R}{\mathbb{R}}
\newcommand{\C}{\mathbb{C}}
\newcommand{\F}{\mathbb{F}}
\newcommand{\A}{\mathbb{A}}
\newcommand{\PP}{\mathbb{P}}
\newcommand{\E}{\mathbb{E}}
\newcommand{\eps}{\varepsilon}
\newcommand{\dd}{\,\mathrm{d}}
\begin{document}
\section*{UnsolvedMath ID: 20002378}
\textbf{Source problem code:} catalogue entry
\subsection{Definitions and mathematical statement}
Source item A.6, ``Dzamonja, Mirna,'' is a participant-interest statement, not a mathematical problem. It describes combinatorial set theory as the participant's primary interest, observes substantial overlap with core model theory in the preceding ten or so years, and expresses an interest in learning core-model methods and understanding their combinatorial nature. The isolated $4$ is a printed page number. No conjecture, quantified target, or request for proof is stated in this source passage.

The following is the supplied research aid's separate classical synthesis, not a question posed by the participant. Work in ZFC, let $L$ be the constructible universe, and let $0^\sharp$ denote zero sharp, the usual sharp for $L$. For a singular strong-limit cardinal $\kappa$, write $\theta=\operatorname{cf}(\kappa)$ and $\lambda=\max\{\aleph_1,\theta\}$. Strong limit means $2^\mu<\kappa$ for every cardinal $\mu<\kappa$. The stated implication is
\[
0^\sharp\text{ does not exist}\quad\Longrightarrow\quad
2^\kappa=\kappa^+\quad\text{for every singular strong-limit }\kappa.
\]
Equivalently, failure of this instance of the singular cardinal hypothesis implies existence of $0^\sharp$; the converse is not asserted.

The aid uses Jensen covering to introduce the family
\[
\mathcal C=\{Y\in\mathcal P(\kappa)\cap L:|Y|^V\leq\lambda\},
\]
where all displayed size bounds and $\kappa^+$, unless marked otherwise, are computed in the ambient universe $V$. Its three reported budgets are
\[
|\mathcal C|^V\leq\kappa^+,\qquad
[\kappa]^\theta\subseteq\bigcup_{Y\in\mathcal C}[Y]^\theta,\qquad
|[Y]^\theta|^V<\kappa\quad(Y\in\mathcal C).
\]
Here $[X]^\theta$ denotes the subsets of $X$ of cardinality $\theta$. In the countable-cofinality case the covering application first enlarges the set being covered to size $\aleph_1$; applying an uncountable-set covering statement directly to a countable set would omit that step. The aid then uses graph coding and K\"onig's theorem to obtain $\kappa^\theta=\kappa^+$, and bounded-initial-segment coding for a strong-limit $\kappa$ to obtain $2^\kappa\leq\kappa^\theta$, yielding the conclusion. These claims are a supplied synthesis of classical inputs, not a new result or a formalization of an actual conjecture in A.6.
\subsection{Short English statement}
A.6 records Mirna Dzamonja's interest in the connection between combinatorial set theory and core model theory. It states no problem. The accompanying aid separately outlines the classical implication from absence of zero sharp to the singular cardinal hypothesis at singular strong-limit cardinals; that synthesis must not be attributed to the participant as an open question.
\subsection{Sources}
\begin{itemize}
\item Local UnsolvedMath registry, permanent ID $20002378$, \texttt{catalogue entry}; \texttt{aim-logic-notes.json}, \texttt{notes[153]}, canonical tag section.
\item American Institute of Mathematics, \emph{Recent advances in core model theory}, Chapter A, ``Participant Contributions,'' A.6, ``Dzamonja, Mirna''; \url{https://aimath.org/WWN/coremodel/coremodel.pdf}. Supplied verification identifies the isolated $4$ as a page-layout artifact, and provides no additional recovered quotation.
\item The research aid names Jensen covering and K\"onig's theorem, without full authoritative bibliographic references for the covering-to-SCH synthesis.
\item Research attempt 1: context only; problem status ``not a problem'' at run; rigor label ``proved''; novelty confidence low. Registry labels open and open triage, and statement status reconstructed/unverified, do not convert the participant statement into a mathematical problem.
\end{itemize}
\end{document}
